\documentclass[10pt, a4paper]{article}
% ===== Essential Packages =====
\usepackage[utf8]{inputenc}       % UTF-8 encoding
\usepackage[T1]{fontenc}          % Better font encoding
\usepackage{lmodern}              % Modern Latin Modern font
\usepackage{ctex}                 % Chinese support (must be loaded after fontenc)
\usepackage{subcaption}
% ===== Math & Symbols =====
\usepackage{amsmath}              % Advanced math
\usepackage{amssymb}              % Additional math symbols
\usepackage{mathrsfs}             % Ralph Smith's Formal Script
\usepackage{braket}               % Dirac bra-ket notation
\usepackage{pdfpages}
% ===== Graphics & Tables =====
\usepackage{graphicx}             % Image support
\usepackage{float}                % Better float control
\usepackage{wrapfig}              % Wrapped figures
\usepackage{tikz}                 % Vector graphics
\usepackage[table]{xcolor}        % Colored tables
\usepackage{array}                % Enhanced tables
\usepackage{multirow}             % Multirow tables
\usepackage{tabularx}             % Adjustable-width tables
\usepackage{booktabs}             % Professional table lines
\usepackage{siunitx}              % SI units in tables
\usepackage{caption}              % Better captions
\usepackage{adjustbox}            % Box resizing

% ===== Code Listings =====
\usepackage{listings}             % Code blocks
\usepackage{tcolorbox}            % Colored boxes
\tcbuselibrary{listings, breakable} % Listings in tcolorbox

% ===== Misc =====
\usepackage{lipsum}               % Dummy text
\usepackage{footmisc}             % Better footnotes
\usepackage[margin=1in]{geometry} % Page margins
\usepackage{diagbox}
\usepackage{ctex}
\author{张世航}
\title{$\gamma$射线的吸收}
\begin{document}
	\maketitle
	\section{各个峰的大致图像}
	\begin{figure}[H]
		\centering
		\includegraphics[width=0.4\linewidth]{newplot(1)}
	\end{figure}
	\begin{tabular}{|c|c|c|}
		\hline
		& CH(道数) & 计数 \\
		\hline
		全能峰右半高& 348 & 372 \\
		\hline
		全能峰左半高& 311 & 371 \\
		\hline
		全能峰& 330 & 737 \\
		\hline
		康普顿边缘& 214 & 220 \\
		\hline
		反散射峰& 54 & 774 \\
		\hline
		X射线峰& 15 & 82 \\
		\hline
	\end{tabular}
	\subsection{定标曲线方程$E = a+b\times CH$}
	我们直接利用最小二乘法进行拟合.我们对这四个值进行拟合.\\
    \begin{tabular}{|c|c|c|c|c|}
    	\hline
    	& 全能峰 & 康普顿边缘& 反散射峰 &  X射线峰\\
    	\hline
    	CH& 330 & 214 & 54 & 15 \\
    	\hline
    	理论能量& 0.662&  0.478& 0.184 & 0.032 \\
    	\hline
    \end{tabular}\\
	\[\includegraphics[width=\textwidth, height=0.2\textheight, keepaspectratio]{2.pdf}\]\\
	在我们处理的结果中,
	\[b = 0.0019306907062004316\]
	\[a = 0.0431216492747839\]
	\[R^2 = 0.9844907334848605\]
	在我们结果中$R^2$很接近1,说明我们的结果大致不错.$b\neq0$说明我们在能量为0的时候道数不为0,可能是因为探测器系统有一个$43.1keV$的能量阈值,可能导致低于这个能量的信号不会被记录或无法被正确分析.\\
	\subsection{能量分辨率}
	我们通过全能峰的数据计算能量分辨率.\\
	\[\eta = \frac{\Delta E}{E}\times100\%\]
	根据我们拟合的结果
	\[E_{\text{全能峰左半高}} = (0.0019306907062004316*311+0.0431216492747839)MeV=0.6435664589031181MeV\]
	\[E_{\text{全能峰右半高}} = (0.0019306907062004316*348+0.0431216492747839)MeV=0.7150020150325341MeV\]
	\[\eta = \frac{E_\text{全能峰右半高}-E_{\text{全能峰左半高}} }{{E_{\text{全能峰}}}} = 10.790869505954076\%\]
	\section{$\gamma$射线的吸收}
	\subsection{Al的吸收}
\begin{tabular}{|c|c|c|c|c|c|c|}
	\hline
	实测厚度& 0 & 9.54mm & 9.44mm & 9.36mm & 9.5mm & 9.46mm \\
	\hline
	计数$n$(个)&167  & 127 & 116 & 105 & 86 & 78 \\
	\hline
	电流强度$I=\frac{n}{t=100s}(A)$&1.67  & 1.27 & 1.16 & 1.05 & 0.86 & 0.78 \\
	\hline
	$lnI$&0.5128 & 0.239 & 0.1484 & 0.04879 & -0.1508 & -0.2485 \\
	\hline
\end{tabular}
\[\includegraphics[width=\textwidth, height=0.2\textheight, keepaspectratio]{Al.pdf}\]\\
\[\mu_{Al} = 0.154cm^{-1}\]

\subsection{Pb的吸收}
\begin{tabular}{|c|c|c|c|c|c|c|}
	\hline
	实测厚度& 0 & 9.54mm & 9.44mm & 9.36mm & 9.5mm & 9.46mm \\
	\hline
	计数$n$(个)&167.0 & 145.0 & 117.0 & 101.0 & 94.0 & 65.0 \\
	\hline
	电流强度$I=\frac{n}{t=100s}(A)$&1.67 & 1.45 & 1.17 & 1.01 & 0.94 & 0.65 \\
	\hline
	$lnI$&0.5128 & 0.3716 & 0.157 & 0.00995 & -0.06188 & -0.4308 \\
	\hline
\end{tabular}
\[\includegraphics[width=\textwidth, height=0.2\textheight, keepaspectratio]{Pb.pdf}\]\\
\[\mu_{Pb} = 1.009cm^{-1}\]
\subsection{Cu的吸收}
\begin{tabular}{|c|c|c|c|c|c|c|}
	\hline
	实测厚度& 0 & 9.54mm & 9.44mm & 9.36mm & 9.5mm & 9.46mm \\
	\hline
	计数$n$(个)&167.0 & 131.0 & 119.0 & 110.0 & 96.0 & 80.0 \\
	\hline
	电流强度$I=\frac{n}{t=100s}(A)$&1.67 & 1.31 & 1.19 & 1.1 & 0.96 & 0.8 \\
	\hline
	$lnI$&0.5128 & 0.27 & 0.174 & 0.09531 & -0.04082 & -0.2231 \\
	\hline
\end{tabular}
\[\includegraphics[width=\textwidth, height=0.2\textheight, keepaspectratio]{Cu.pdf}\]\\
\[\mu_{Cu} = 0.54cm^{-1}\]


\end{document} 